\subsection{Point-source obstacle form and semantic target}
\label{subsec:two-stage-obstacle}

For the remainder of the note let $s=e_v$.  Put
\begin{equation}
 c_\rho:=b-\alpha\rho D^{1/2}\one.
 \label{eq:two-stage-load}
\end{equation}
The source nonnegativity theorem converts RPPR into the Stieltjes obstacle
problem
\begin{equation}
 x^\star(\rho)
 =\argmin_{x\geq0}
 \left\{\frac12\langle x,Qx\rangle-\langle c_\rho,x\rangle\right\}.
 \label{eq:two-stage-obstacle}
\end{equation}
For $x\geq0$, define the boundary key
\begin{equation}
 g_i(x):=(c_\rho-Qx)_i.
 \label{eq:two-stage-key}
\end{equation}
At the optimum, $g_i(x^\star)=0$ on $S^\star$ and
$g_i(x^\star)\leq0$ off $S^\star$.

The note-scoped PPR target is
\begin{equation}
 \norm{D^{-1/2}(\widehat x-x^0)}_\infty\leq\epsppr.
 \label{eq:two-stage-ppr-target}
\end{equation}
We use the source comparison
\begin{equation}
 0\leq D^{-1/2}(x^0-x^\star(\rho))\leq\rho\one,
 \qquad
 \vol(S^\star(\rho))\leq\frac1\rho.
 \label{eq:two-stage-rppr-bridge}
\end{equation}

\begin{lemma}[The point-source active face has one root]
\label{lem:two-stage-point-source-connected}
If $S^\star(\rho)$ is nonempty, it contains the seed and induces a connected
subgraph.
\end{lemma}

\begin{proof}
Let $C$ be a connected component of the active induced subgraph that does not
contain the seed.  Active KKT equalities and the absence of active edges from
$C$ to the other components give
\begin{equation*}
 Q_Cx_C^\star=c_{\rho,C}
 =-\alpha\rho D_C^{1/2}\one<0.
\end{equation*}
The principal Stieltjes inverse is nonnegative, so this would imply
$x_C^\star<0$, contradicting activity.  Therefore every active component
contains the unique seed; there can be at most one.
\end{proof}

One sufficient route is to set
\begin{equation}
 \rho=\tau=\frac{\epsppr}{2}
 \label{eq:two-stage-split}
\end{equation}
and return an RPPR approximation satisfying
$\norm{D^{-1/2}(\widehat x-x^\star(\rho))}_\infty\leq\tau$.
The stronger envelope-linearization result below gives a simpler second
route: after using RPPR only to certify a set, solve ordinary PPR on that set.

\subsection{Exact-face and containing-envelope certificates}

\begin{definition}[Exact-face certificate]
\label{def:two-stage-face-certificate}
An exact-face certificate is either the certified empty optimum, or a set
$A\ni v$ and a vector $p_A>0$ such that
\begin{equation}
 Q_Ap_A=c_{\rho,A},
 \qquad
 c_{\rho,j}-Q_{jA}p_A\leq0\quad(j\in\partial A).
 \label{eq:two-stage-face-certificate}
\end{equation}
It may additionally contain a lower checkpoint $0\leq z\leq p_A$ satisfying
$c_{\rho,A}-Q_Az\geq0$.
\end{definition}

Because every nonseed load outside $A\cup\partial A$ is strictly negative
and has no edge into $A$, the two conditions in
\cref{eq:two-stage-face-certificate}, together with zero padding, are the
global obstacle KKT conditions.  Hence $A=S^\star$ and $p_A=x_A^\star$.

\begin{definition}[Containing-envelope certificate]
\label{def:two-stage-envelope-certificate}
A containing-envelope certificate is a finite set $E$ for which the
producing algorithm proves
\begin{equation}
 S^\star(\rho)\subseteq E,
 \qquad
 \vol(E)\leq V_E.
 \label{eq:two-stage-envelope-certificate}
\end{equation}
The support inclusion may be a theorem of the producing algorithm; it need
not reveal the exact active face within $E$.
\end{definition}

The second certificate is deliberately weaker.  Once
$S^\star\subseteq E$, restricting the obstacle problem to $E$ does not change
its minimizer.  This observation removes exact face identification, strict
complementarity, and zero-key decisions from the minimal end-to-end theorem.

\begin{definition}[Approximate-envelope certificate]
\label{def:two-stage-approximate-envelope}
For $\delta\geq0$, an approximate-envelope certificate is a finite set $E$
and a proof that
\begin{equation}
 \max_{i\notin E}\frac{x_i^\star(\rho)}{\sqrt{d_i}}\leq\delta.
 \label{eq:two-stage-approximate-envelope}
\end{equation}
A containing-envelope certificate is the special case $\delta=0$.
\end{definition}

\begin{definition}[Mass-completion certificate]
\label{def:two-stage-mass-certificate}
A mass-completion certificate is a sparse vector $z$ and a declared constant
$c>0$ such that
\begin{equation}
 0\leq z,\qquad b-Qz\geq0,\qquad
 1-\langle D^{1/2}\one,z\rangle\leq c\sqrt\alpha.
 \label{eq:two-stage-mass-certificate}
\end{equation}
It is a numerical certificate, but it need not yet meet the requested PPR
accuracy and it need not contain, approximate, or identify the RPPR support.
The remaining positive residual mass is small enough for a monotone
APPR/SOR cleanup.
\end{definition}

Thus Stage~I has three useful exits: a target-accurate numerical point, a
set-only envelope, or a lower-safe mass-completion point.  Only the set-only
exit invokes a new fixed principal linear solve.  The mass exit continues
from $z$ and is analyzed in
\cref{prop:two-stage-aesp-appr-mass-handoff,thm:two-stage-mass-capturing-envelope}.

There is also a direct set certificate that bypasses RPPR altogether.  It is
especially useful when a candidate envelope comes from graph structure rather
than an obstacle solve.  Put $a=(1+\alpha)/2$ and
$c=(1-\alpha)/2$.

\begin{definition}[Observable PPR boundary-leakage certificate]
\label{def:two-stage-leakage-certificate}
Let $E$ contain the point source, let $u_E=Q_E^{-1}b_E$, and suppose a
computed vector $\widehat u_E$ has the certified error bar
\begin{equation}
 \max_{i\in E}
 \frac{|(u_E)_i-(\widehat u_E)_i|}{\sqrt{d_i}}
 \leq\eta.
 \label{eq:two-stage-leakage-error-bar}
\end{equation}
For $j\in\partial E$, define the observable leakage estimate
\begin{equation}
 \widehat\ell_j
 :=\frac{-Q_{jE}\widehat u_E}{\sqrt{d_j}}
 =\frac{c}{d_j}
   \sum_{i\in E\cap N(j)}
   \frac{(\widehat u_E)_i}{\sqrt{d_i}}.
 \label{eq:two-stage-observed-leakage}
\end{equation}
The pair $(E,\widehat u_E)$ is a $\delta$ boundary-leakage certificate if
\begin{equation}
 \widehat\ell_j
 +c\frac{\deg_E(j)}{d_j}\eta
 \leq\alpha\delta
 \qquad(j\in\partial E).
 \label{eq:two-stage-leakage-gate}
\end{equation}
\end{definition}

\begin{proposition}[A local leakage gate certifies a PPR envelope]
\label{prop:two-stage-leakage-certificate}
If \cref{eq:two-stage-leakage-gate} holds, then the zero extension
$\widetilde u_E$ satisfies
\begin{equation}
 0\leq\widetilde u_E\leq x^0,
 \qquad
 \norm{D^{-1/2}(x^0-\widetilde u_E)}_\infty\leq\delta.
 \label{eq:two-stage-leakage-conclusion}
\end{equation}
Consequently a fixed-envelope Stage~II approximation with normalized error
at most $\tau$ is a full-PPR $(\delta+\tau)$-approximation.  A residual-based
fully observable choice is
\begin{equation}
 \eta=\frac{\norm{b_E-Q_E\widehat u_E}_2}{\alpha},
 \label{eq:two-stage-leakage-residual-error}
\end{equation}
because $Q_E\succeq\alpha I$ and $d_i\geq1$.
\end{proposition}

\begin{proof}
Pad $u_E$ by zero and put $r=b-Q\widetilde u_E$.  Its rows in $E$ vanish.
Outside $E$ the point-source load vanishes, so
\begin{equation*}
 r_j=-Q_{jE}u_E\geq0,
\end{equation*}
and nonboundary rows vanish as well.  The error bar gives
\begin{equation*}
 \frac{r_j}{\sqrt{d_j}}
 \leq\widehat\ell_j
 +c\frac{\deg_E(j)}{d_j}\eta
 \leq\alpha\delta.
\end{equation*}
Thus $0\leq r\leq\alpha\delta D^{1/2}\one$.  Since
$QD^{1/2}\one=\alpha D^{1/2}\one$ and $Q^{-1}\geq0$,
\begin{equation*}
 0\leq x^0-\widetilde u_E=Q^{-1}r
 \leq\delta D^{1/2}\one.
\end{equation*}
This proves \cref{eq:two-stage-leakage-conclusion}; the terminal error follows
by the triangle inequality.  Finally, strong convexity gives
$\norm{u_E-\widehat u_E}_2\leq
\norm{b_E-Q_E\widehat u_E}_2/\alpha$, which implies
\cref{eq:two-stage-leakage-error-bar} with the declared $\eta$.
\end{proof}

\begin{corollary}[An APPR/SOR checkpoint screens the fixed envelope]
\label{cor:two-stage-monotone-leakage-screen}
Let $z\geq0$ be a global PPR subsolution, $r=b-Qz\geq0$, and
$E=\supp(z)\ni v$.  Put
\begin{equation}
 \eta_E:=\max_{i\in E}\frac{r_i}{\sqrt{d_i}}.
 \label{eq:two-stage-internal-key-scale}
\end{equation}
Then the exact principal solution obeys
\begin{equation}
 0\leq u_E-z_E
 \leq\frac{\eta_E}{\alpha}D_E^{1/2}\one.
 \label{eq:two-stage-monotone-principal-box}
\end{equation}
Hence the fully observable gate
\begin{equation}
 \frac{r_j}{\sqrt{d_j}}
 +c\frac{\deg_E(j)}{d_j}\frac{\eta_E}{\alpha}
 \leq\alpha\delta
 \qquad(j\in\partial E)
 \label{eq:two-stage-monotone-leakage-gate}
\end{equation}
is sufficient for a $\delta$ boundary-leakage certificate.  If it passes
before the global APPR stopping test, discarding the monotone history and
accelerating $Q_Eu_E=b_E$ gives a genuine two-stage continuation rather than
a redundant tail.
The still cheaper two-heap sufficient test is
\begin{equation}
 \eta_{\partial E}+\frac{c}{\alpha}\eta_E
 \leq\alpha\delta,
 \qquad
 \eta_{\partial E}:=
 \max_{j\in\partial E}\frac{r_j}{\sqrt{d_j}},
 \label{eq:two-stage-two-heap-leakage-gate}
\end{equation}
with an empty-boundary maximum equal to zero.  It follows from
$\deg_E(j)/d_j\leq1$ and can be maintained by one lazy heap for internal rows
and one for boundary rows, without a boundary sweep.  When
$\partial E=\varnothing$, the leakage gate is vacuous; the displayed
two-term test is used only for a nonempty boundary.
\end{corollary}

\begin{proof}
The principal error satisfies
$Q_E(u_E-z_E)=r_E\leq\eta_ED_E^{1/2}\one$.  Since
$Q_ED_E^{1/2}\one\geq\alpha D_E^{1/2}\one$, inverse positivity proves
\cref{eq:two-stage-monotone-principal-box}.  The current exterior residual
density is exactly the leakage estimate
$r_j/\sqrt{d_j}=-Q_{jE}z_E/\sqrt{d_j}$.  Substituting
$\eta=\eta_E/\alpha$ in
\cref{eq:two-stage-leakage-gate} proves the claim.
\end{proof}

This gate is the shortest literal two-stage interface in the note.  Stage~I
may propose $E$ by APPR, SOR, AESP-CD, a rooted ball, or a structural
reporter; a coarse fixed-envelope solve and one boundary scan decide whether
the set is sufficient.  If it passes, Stage~II refines the same principal
system once.  No RPPR support, zero-key decision, or obstacle history is part
of the certificate.  If it fails, the positive leakage rows name where the
envelope must grow, subject to the same hard cap and fallback discipline as
the other lanes.

The following further reduction is the key two-stage simplification.  Let
$u_E$ solve the \emph{unregularized} principal PPR system
\begin{equation}
 Q_Eu_E=b_E,
 \qquad \widetilde u_E:=\text{the zero extension of }u_E.
 \label{eq:two-stage-restricted-ppr}
\end{equation}

\begin{proposition}[An approximate RPPR envelope linearizes Stage~II]
\label{prop:two-stage-envelope-linearization}
For any $E$, put
$\delta_E:=\max_{i\notin E}x_i^\star(\rho)/\sqrt{d_i}$, with the maximum of
the empty set equal to zero.  Then
\begin{equation}
 0\leq u_E\leq x_E^0,
 \qquad
 \norm{D^{-1/2}(x^0-\widetilde u_E)}_\infty
 \leq\rho+\delta_E.
 \label{eq:two-stage-envelope-linearization}
\end{equation}
In particular, a $\delta$-approximate-envelope certificate incurs at most
$\rho+\delta$ truncation error.  A containing envelope incurs at most
$\rho$.  Stage~II may therefore solve
\cref{eq:two-stage-restricted-ppr}; it need not solve the obstacle problem
again or classify the coordinates inside $E$.
\end{proposition}

\begin{proof}
The RPPR comparison gives, outside $E$,
\begin{equation*}
 0\leq x_i^0\leq x_i^\star(\rho)+\rho\sqrt{d_i}
 \leq(\rho+\delta_E)\sqrt{d_i}.
\end{equation*}
Put $w=x_E^0-u_E$.  The full and principal PPR equations give
\begin{equation*}
 Q_Ew=-Q_{E\bar E}x_{\bar E}^0\geq0,
\end{equation*}
so inverse positivity gives $0\leq u_E\leq x_E^0$.  Let
$\gamma=\rho+\delta_E$.  Since
$x_{\bar E}^0\leq\gamma D_{\bar E}^{1/2}\one$ and
$Q_{E\bar E}\leq0$,
\begin{equation*}
 Q_Ew
 \leq-Q_{E\bar E}\gamma D_{\bar E}^{1/2}\one
 \leq Q_E\gamma D_E^{1/2}\one.
\end{equation*}
The last inequality is the restriction of
$QD^{1/2}\one=\alpha D^{1/2}\one$.  A final use of principal inverse
positivity yields $0\leq w\leq\gamma D_E^{1/2}\one$.  Together with the
outside bound, this proves \cref{eq:two-stage-envelope-linearization}.
\end{proof}

\begin{remark}[The linearization is not point-source specific]
The proof uses only $b\geq0$, the Stieltjes signs, and an envelope containing
the RPPR support.  It therefore extends to every nonnegative source
distribution.  The point-source restriction is needed for the strongest
known Stage~I locality, connected-support, and rooted-reporter results, not
for the fixed-envelope linear solve.
\end{remark}

\subsection{Charged work}

Scanning row $i$ costs $d_i$.  Stage~I work includes every exposure, repeated
row read, coordinate or batch update, boundary query, response or factor
update, certificate scan, and state write.  Stage~II charges every restricted
linear-system matrix application and final output write.  Preprocessing
outside the exposed graph is not free.  All exact-real structural results and finite-precision
claims are kept separate.
